% ============================================================================
\chapter{Калибровочная группа Стандартной модели, гравитационное преломление и космология}
\label{ch:mono-sm-gravity}
% ============================================================================

\section{Симплициальный комплекс \texorpdfstring{$\mathcal{C}_{13}$}{C13}, матричные коммутаторы алгебры Ли \texorpdfstring{$\mathfrak{u}(1)\oplus\mathfrak{su}(2)\oplus\mathfrak{su}(3)$}{u(1) + su(2) + su(3)} и четырёхпетлевой вывод постоянной тонкой структуры}

\begin{theorem}[Кликовый комплекс $\mathcal{C}_{13}$, матричные коммутаторы Ли $\mathfrak{u}(1)\oplus\mathfrak{su}(2)\oplus\mathfrak{su}(3)$ и четырёхпетлевой вывод $\alpha_{\mathrm{fs}}^{-1}(3) \in \mathbb{Q}(\sqrt{2})$]
\label{thm:mono-c13-lie-alpha-fs}
На $13$-узловом причинном ядре кубооктаэдра $B_1(3) = \{\mathbf{0}\} \cup O_2(3)$ выполняются следующие точные структурные законы (\texttt{sm\_\allowbreak and\_\allowbreak cosmology\_\allowbreak closed\_\allowbreak rational\_\allowbreak operators}, \texttt{gaussian\_\allowbreak coulomb\_\allowbreak lock\_\allowbreak iff\_\allowbreak d3}, \texttt{alpha\_\allowbreak fs\_\allowbreak exact\_\allowbreak z\_\allowbreak sqrt2\_\allowbreak canonical\_\allowbreak form\_\allowbreak and\_\allowbreak floor\_\allowbreak 137} в модуле \texttt{lean/MtCarrier/ContinuumPhysics.lean}):
\begin{enumerate}[label=\textbf{(\arabic*)},leftmargin=*]
\item \textbf{Перепись клик $\mathcal{C}_{13}$ и эйлерова характеристика $\chi(\mathcal{C}_{13}) = 1$:}
  числа полных подграфов (клик $K_n$) в $B_1(3)$ и знакопеременная эйлерова характеристика равны:
  \begin{align}
    \bigl(K_1(3),\ K_2(3),\ K_3(3),\ K_4(3),\ K_{\ge 5}(3)\bigr) &= (13,\ 36,\ 32,\ 8,\ 0), \label{eq:mono-c13-cliques} \\
    \chi(\mathcal{C}_{13}) \;=\; K_1 - K_2 + K_3 - K_4 &= 13 - 36 + 32 - 8 \;=\; 1 \;=\; \mu_{\mathrm{vac}}(3). \label{eq:mono-c13-euler}
  \end{align}
\item \textbf{Явные целочисленные матричные коммутаторы $[A, B] = AB - BA$ алгебры Ли $\mathfrak{u}(1) \oplus \mathfrak{su}(2) \oplus \mathfrak{su}(3)$ ($\dim = 1 + 3 + 8 = 12 = |N(3)|$):}
  пространство независимых циклов голономии расщепляется в прямую сумму
  \begin{equation}
    \label{eq:mono-sm-gauge}
    \dim\bigl(\mathfrak{u}(1) \oplus \mathfrak{su}(2) \oplus \mathfrak{su}(3)\bigr) \;=\; 1 + 3 + 8 \;=\; 12 \;=\; |N(3)|,
  \end{equation}
  где: (а)~$1$ центральный фазовый генератор $Y_0 = I_3$ на $\mathbb{Z}_{512}$ коммутирует со всеми генераторами ($[Y_0, L_k] = 0$, $[Y_0, E_{ab}] = 0$) и порождает $\mathfrak{u}(1)_Y$; (б)~$d = 3$ антисимметричные матрицы инфинитезимальных поворотов $(L_k)_{ij} = -\varepsilon_{kij}$ в трёх координатных плоскостях $(yz, zx, xy)$, пересекающихся исключительно в центральном узле $\mathbf{0}$, удовлетворяют коммутационным соотношениям $\mathfrak{su}(2)$:
  \begin{equation}
    \label{eq:mono-su2-comm}
    [L_1, L_2] \;=\; L_3, \qquad
    [L_2, L_3] \;=\; L_1, \qquad
    [L_3, L_1] \;=\; L_2;
  \end{equation}
  (в)~на $K_4(3) = 8$ тетраэдрических 4-кликах три цветовые плоскости каждого радиального бариона $K_3$ порождают $3^2 - 1 = 8$ бесследовых матричных генераторов $\mathfrak{su}(3)$ ($6$ корневых переходов $E_{ab}$ при $a \ne b \in \{1,2,3\}$ и $2$ диагональных генератора Картана $H_1 = \operatorname{diag}(1,-1,0)$, $H_2 = \operatorname{diag}(0,1,-1)$), удовлетворяющих соотношениям Шевалле--Серра:
  \begin{align}
    [E_{12}, E_{23}] &= E_{13}, \qquad
    [E_{12}, E_{21}] \;=\; H_1, \qquad
    [E_{23}, E_{32}] \;=\; H_2, \label{eq:mono-su3-comm} \\
    [H_1, H_2] &= 0, \qquad
    \operatorname{Tr}(H_1) \;=\; \operatorname{Tr}(H_2) \;=\; 0. \label{eq:mono-su3-cartan}
  \end{align}
\item \textbf{Единственность гауссова кулоновского модуля $\widehat{M}(3) = 97$, точный обрыв симплициального ряда на $k_{\max} = 4$ ($K_{\ge 5}(\mathcal{C}_{13}) \equiv 0$) и вывод $\alpha_{\mathrm{fs}}^{-1}(3) \in \mathbb{Q}(\sqrt{2})$:}
  полное число направленных вершин-октантов причинного ядра $1 + |O_2(d)|\cdot|{\triangle}(d)|$ совпадает с квадратом нормы гауссова фильтра $|d_{\mathrm{in}}(\square) + i\,\mathrm{mod\_bits}(d)|^2$ в $\mathbb{Z}[i]$ тогда и только тогда, когда $d = 3$ (\texttt{gaussian\_\allowbreak coulomb\_\allowbreak lock\_\allowbreak iff\_\allowbreak d3}), определяя единственный простой гауссов кулоновский модуль $\widehat{M}(3)$ и целую часть обратной постоянной тонкой структуры:
  \begin{align}
    \widehat{M}(3) &= 1 + |O_2(3)|\cdot|{\triangle}(3)| \;=\; d_{\mathrm{in}}(\square)^2 + \mathrm{mod\_bits}(3)^2 \;=\; 4^2 + 9^2 \;=\; 97, \label{eq:mono-mhat-97} \\
    \bigl\lfloor \alpha_{\mathrm{fs}}^{-1}(3) \bigr\rfloor &= \widehat{M}(3) + \bigl(\Delta\varphi_{\mathrm{disc}}(3) - 1\bigr) \;=\; 97 + 40 \;=\; 137 \;=\; Q_\square(3) + \mathrm{mod\_bits}(3). \label{eq:mono-alpha-fs}
  \end{align}
  Поскольку постоянная тонкой структуры $\alpha_{\mathrm{fs}}(d)$ является статическим геометрическим инвариантом самой координационной решётки причинного ядра $\mathcal{C}_{13}$ (не зависящим от такта времени $t$ и конкретной динамической траектории поля), а по пункту~(1)~\eqref{eq:mono-c13-cliques} кликовый комплекс причинного ядра $\mathcal{C}_{13}$ исчерпывается симплексами размерностей $k \in \{1, 2, 3, 4\}$ ($K_1 \subset K_2 \subset K_3 \subset K_4$) и \textbf{не содержит симплексов порядка $k \ge 5$ ($K_{\ge 5}(\mathcal{C}_{13}) \equiv 0$)}, разложение статического кулоновского пропагатора по невырожденным симплициальным путям внутри $\mathcal{C}_{13}$ с кулоновским весом $\widehat{M}(d)^{-1}$ на каждую вершину симплекса \textbf{обрывается строго на четвёртом порядке $k_{\max} = 4$} (вклады $\widehat{M}(d)^{-k}$ при всех $k \ge 5$ тождественно равны нулю). На четырёх непустых этажах симплициальной фильтрации $K_1 \to K_2 \to K_3 \to K_4$ вдоль диагонального ребра $\|\mathbf{e}\|_2 = \sqrt{d-1} = \sqrt{2}$ с осевым проекционным фактором $\kappa_{\mathrm{geom}} = 1/\sqrt{2} = \sqrt{2}/2$, числом открытых внешних октантов вершины $d_{\mathrm{out}}(V) = |{\triangle}(d)| - 1 = 7$ и квадратным входом $d_{\mathrm{in}}(\square) = 4$ числитель вакуумной экранировки $\mathcal{N}_\alpha(d)$ ($k = 3, 4$) и знаменатель прямой проводимости звена $\mathcal{D}_\alpha(d)$ ($k = 1, 2$) однозначно определяются комбинаторными весами соответствующих симплексов:
  \begin{align}
    \mathcal{D}_\alpha(d) &= \underbrace{\frac{\kappa_{\mathrm{geom}}}{\widehat{M}(d)}}_{k=1\ (K_2\text{-ребро})} \;+\; \underbrace{\frac{\kappa_{\mathrm{geom}}^2}{d_{\mathrm{out}}(V)\,\widehat{M}(d)^2}}_{k=2\ (K_3\text{-возвратный клин})} \;=\; \frac{\sqrt{2}}{2\,\widehat{M}(d)} + \frac{1}{2\bigl(|{\triangle}(d)|-1\bigr)\widehat{M}(d)^2}, \label{eq:mono-alpha-4loop-den} \\
    \mathcal{N}_\alpha(d) &= 1 \;-\; \underbrace{\frac{2\kappa_{\mathrm{geom}}}{d_{\mathrm{in}}(\square)\,d_{\mathrm{out}}(V)\,\widehat{M}(d)^3}}_{k=3\ (\square\to K_4\text{ экранировка грани})} \;-\; \underbrace{\frac{|K_4(d)|}{d_{\mathrm{out}}(V)^2\,\widehat{M}(d)^4}}_{k=4\ (\text{замыкание }K_4)} \nonumber \\
    &= 1 - \frac{\sqrt{2}}{d_{\mathrm{in}}(\square)\bigl(|{\triangle}(d)|-1\bigr)\widehat{M}(d)^3} - \frac{|{\triangle}(d)|}{\bigl(|{\triangle}(d)|-1\bigr)^2\widehat{M}(d)^4}, \label{eq:mono-alpha-4loop-num} \\
    \alpha_{\mathrm{fs}}^{-1}(d) &= \frac{\mathcal{N}_\alpha(d)}{\mathcal{D}_\alpha(d)}. \label{eq:mono-alpha-4loop-fraction}
  \end{align}
  Приведение~\eqref{eq:mono-alpha-4loop-den}--\eqref{eq:mono-alpha-4loop-fraction} к общему целому знаменателю $d_{\mathrm{in}}(\square)\bigl(|{\triangle}(d)|-1\bigr)^2\widehat{M}(d)^4$ определяет коэффициенты в кольце $\mathbb{Z}[\sqrt{2}]$:
  \begin{align}
    A_{\mathrm{num}}(d) &= \tfrac{1}{2}d_{\mathrm{in}}(\square)\bigl(|{\triangle}(d)|-1\bigr)\hat{M}(d)^2 \;=\; 131\,726, \label{eq:mono-alpha-zsqrt2-coeffs} \\
    B_{\mathrm{num}}(d) &= \tfrac{1}{2}d_{\mathrm{in}}(\square)\bigl(|{\triangle}(d)|-1\bigr)^2\hat{M}(d)^3 \;=\; 89\,441\,954, \label{eq:mono-alpha-bnum} \\
    A_{\mathrm{den}}(d) &= d_{\mathrm{in}}(\square)\Bigl(\bigl(|{\triangle}(d)|-1\bigr)^2\hat{M}(d)^4 - |{\triangle}(d)|\Bigr) \;=\; 17\,351\,739\,044, \label{eq:mono-alpha-aden} \\
    B_{\mathrm{den}}(d) &= -\bigl(|{\triangle}(d)|-1\bigr)\hat{M}(d) \;=\; -679, \label{eq:mono-alpha-bden}
  \end{align}
  и замкнутую алгебраическую формулу в поле $\mathbb{Q}(\sqrt{2})$ без трансцендентного числа~$\pi$:
  \begin{equation}
    \label{eq:mono-alpha-codata}
    \alpha_{\mathrm{fs}}^{-1}(3) \;=\; \frac{A_{\mathrm{den}}(3) + B_{\mathrm{den}}(3)\sqrt{2}}{A_{\mathrm{num}}(3) + B_{\mathrm{num}}(3)\sqrt{2}} \;=\; \frac{17\,351\,739\,044 - 679\sqrt{2}}{131\,726 + 89\,441\,954\sqrt{2}} \;=\; 137{,}035\,999\,187\dots
  \end{equation}
  (CODATA-2022: $137{,}035\,999\,177(21)$, отклонение $0{,}47$ стандартного отклонения).
\end{enumerate}
\end{theorem}

\begin{proof}
В пункте~(1) перепись полных подграфов $B_1(3)$ даёт $13$ вершин ($K_1 = 13$), $12$ радиальных и $24$ поверхностных ребра ($K_2 = 36$), $24$ радиальных и $8$ поверхностных треугольников ($K_3 = 32$) и $8$ тетраэдров ($K_4 = 8$). Поскольку любые четыре взаимно смежные вершины первой сферы $O_2(3)$ с центром $\mathbf{0}$ образуют правильный тетраэдр $K_4$ в одном из $8$ октантов, а пятая вершина из $O_2(3)$ находится на квадрате расстояния $\ge 4 > 2$ хотя бы от одной из вершин этого тетраэдра, в $B_1(3)$ не существует ни одного полного подграфа из $5$ и более вершин ($K_{\ge 5}(3) \equiv 0$), откуда $\chi(\mathcal{C}_{13}) = 13 - 36 + 32 - 8 = 1$.

В пункте~(2) прямое перемножение целочисленных $3\times 3$ матриц $L_1, L_2, L_3$ и $E_{ab}, H_1, H_2$ проверяет коммутаторы~\eqref{eq:mono-su2-comm}--\eqref{eq:mono-su3-cartan}.

В пункте~(3) равенство $1 + 2d(d-1)\cdot 2^d = (2(d-1))^2 + (d^2)^2$ при $d = 3$ даёт $1 + 12\times 8 = 97 = 4^2 + 9^2 = \widehat{M}(3)$ (тогда как при $d \ne 3$ левая и правая части различны). Так как по пункту~(1) $K_{\ge 5}(\mathcal{C}_{13}) \equiv 0$, любой невырожденный симплициальный путь по цепочке вложенных клик внутри причинного ядра $B_1(3)$ содержит от $k = 1$ до $k_{\max} = 4$ шагов вершинного подавления $\widehat{M}(3)^{-1}$, поэтому члены порядка $\widehat{M}(3)^{-5}$ и выше на комплексе $\mathcal{C}_{13}$ отсутствуют тождественно. Вычислим комбинаторный вес каждого из четырёх симплициальных порядков $k \in \{1, 2, 3, 4\}$:
\begin{itemize}[leftmargin=*,noitemsep]
\item \emph{Порядок $k = 1$ (прямой перенос вдоль ребра $K_2$ в знаменателе проводимости $\mathcal{D}_\alpha$):} единичный шаг вдоль диагонального ребра $\mathbf{e} \in O_2(3)$ длины $\|\mathbf{e}\|_2 = \sqrt{2}$ проектируется на координатную ось с множителем $\cos 45^\circ = \kappa_{\mathrm{geom}} = 1/\sqrt{2} = \sqrt{2}/2$, давая вклад $+\frac{\sqrt{2}}{2}\widehat{M}(d)^{-1}$;
\item \emph{Порядок $k = 2$ (однократное возвратное рассеяние на 2-симплексе $K_3$ в $\mathcal{D}_\alpha$):} возвратный переход по двум рёбрам треугольной грани имеет проекционный квадрат $\kappa_{\mathrm{geom}}^2 = 1/2$ и равномерно распределяется между $d_{\mathrm{out}}(V) = |{\triangle}(d)| - 1 = 7$ открытыми внешними октантами вершины оболочки, давая вклад $+\frac{\kappa_{\mathrm{geom}}^2}{d_{\mathrm{out}}(V)}\widehat{M}(d)^{-2} = +\frac{1}{2\cdot 7}\widehat{M}(d)^{-2}$;
\item \emph{Порядок $k = 3$ (замкнутая экранирующая петля по квадратной грани $\square$ с выходом в открытые октанты в числителе $\mathcal{N}_\alpha$):} трёхзвенное замыкание вокруг квадратной грани из $d_{\mathrm{in}}(\square) = 4$ вершин с двумя противоположными парами рёбер ($2\kappa_{\mathrm{geom}} = \sqrt{2}$), нормированное на число пар входа-выхода $d_{\mathrm{in}}(\square)\,d_{\mathrm{out}}(V) = 4 \times 7 = 28$, вычитает из единицы вакуумную экранировку третьего порядка $-\frac{\sqrt{2}}{d_{\mathrm{in}}(\square)(|{\triangle}(d)|-1)}\widehat{M}(d)^{-3}$;
\item \emph{Порядок $k = 4$ (терминальное замыкание всех тетраэдров $K_4$ в числителе $\mathcal{N}_\alpha$):} полный обход всех $|K_4(d)| = |{\triangle}(d)| = 8$ четырёхвершинных тетраэдров ядра, нормированный на квадрат числа открытых внешних октантов пары граничных вершин $d_{\mathrm{out}}(V)^2 = (|{\triangle}(d)| - 1)^2 = 49$, вычитает терминальный вклад четвёртого порядка $-\frac{|{\triangle}(d)|}{(|{\triangle}(d)|-1)^2}\widehat{M}(d)^{-4}$.
\end{itemize}
Умножение числителя и знаменателя~\eqref{eq:mono-alpha-4loop-fraction} на общий целый знаменатель $d_{\mathrm{in}}(\square)\bigl(|{\triangle}(d)|-1\bigr)^2\widehat{M}(d)^4 = 4 \times 49 \times 97^4$ даёт коэффициенты~\eqref{eq:mono-alpha-zsqrt2-coeffs}--\eqref{eq:mono-alpha-bden}, а избавление от иррациональности в знаменателе приводит дробь к каноническому виду~\eqref{eq:mono-alpha-codata} в $\mathbb{Q}(\sqrt{2})$ с целой частью $137$.
\end{proof}

\section{Структура 13-узлового ротора электрона, полное квантование восьми наблюдаемых и единый вихревой атлас Стандартной модели на кликах \texorpdfstring{$\mathcal{C}_{13}$}{C13}}

\begin{theorem}[Этаж~1 и полное квантование: 13-узловой одетый ротор электрона, расслоение Вейля--Швингера и квантование всех восьми физических наблюдаемых]
\label{thm:mono-electron-rotor-quantization}
На причинном ядре $B_1(3) = \{\mathbf{0}\}\cup O_2(3)$ с однотактовым спинорным бюджетом ячейки $B_1(3) = d_{\mathrm{in}}(\square)\,\mathrm{mod\_bits}(3) = 4\times 9 = 36$ бит и двухтактовым динамическим бюджетом фазового пространства $B_2(3) = 2 B_1(3) = 72$ бита выполняются следующие законы (\texttt{complete\_\allowbreak quantization\_\allowbreak of\_\allowbreak phase\_\allowbreak space\_\allowbreak and\_\allowbreak all\_\allowbreak observables\_\allowbreak operator}, \texttt{elementary\_\allowbreak particles\_\allowbreak electron\_\allowbreak rotor\_\allowbreak clique\_\allowbreak complex\_\allowbreak and\_\allowbreak nucleon\_\allowbreak moments\_\allowbreak operator} в модуле \texttt{lean/MtCarrier/ParticleMaterials.lean}):
\begin{enumerate}[label=\textbf{(\arabic*)},leftmargin=*]
\item \textbf{Минимальный радиус $\varepsilon$-ореола $R_{\mathrm{dress}} = 1\cdot dh$ и разбиение $13$-узлового ротора электрона $(1 + 1 + 4 + 7 = 13)$:}
  поскольку фазовый шаг за полный оборот по $12$ соседям первой сферы превышает порог дискретности ($\lfloor 512/12 \rfloor = 42 > 41 = \Delta\varphi_{\mathrm{disc}}(3)$), элементарный единичный заряд $Q = \pm 1$ не может быть локализован в одной изолированной вершине и неизбежно возбуждает всю первую координационную сферу $O_2(3)$ из $12$ соседей ($R_{\mathrm{dress}} = 1\cdot dh$). Относительно любой опорной вершины оболочки $v_0 = (0, 1, 1) \in O_2(3)$ все $13$ узлов одетого кластера электрона $B_1(3)$ разбиваются по квадрату евклидова расстояния $|u - v_0|_2^2$ на четыре слоя:
  \begin{equation}
    \label{eq:mono-electron-13-partition}
    N_\varphi(3) \;=\; 1_{\mathrm{core}} \;+\; d_{\mathrm{in}}(V) \;+\; d_{\mathrm{same}}(V) \;+\; d_{\mathrm{out}}(V) \;=\; 1 + 1 + 4 + 7 \;=\; 13,
  \end{equation}
  где $1_{\mathrm{core}} = 1$ --- центральный узел $(0,0,0)$, $d_{\mathrm{in}}(V) = 1$ --- сама опорная вершина $v_0$ ($|u - v_0|_2^2 = 0$), $d_{\mathrm{same}}(V) = d_{\mathrm{in}}(\square) = 4$ --- четыре смежные вершины оболочки на расстоянии $|u - v_0|_2^2 = 2$, а $d_{\mathrm{out}}(V) = |{\triangle}(3)| - 1 = 7$ --- семь несмежных вершин оболочки на расстоянии $|u - v_0|_2^2 \in \{4, 6, 8\}$.
\item \textbf{Расслоение Вейля--Швингера и полное квантование всех восьми физических наблюдаемых:}
  каждая фаза кольца $\varphi \in \{0, \dots, 511\}$ взаимно-однозначно раскладывается по формуле $\varphi = 64 R + q$ на номер октантного уровня $R = \lfloor\varphi / 64\rfloor \in \mathbb{Z}_8$ и мантиссу $q = \varphi \bmod 64 \in \mathbb{Z}_{64}$ ($\mathbb{Z}_{512} \cong \mathbb{Z}_8 \times \mathbb{Z}_{64}$). На $M$-носителе все восемь физических наблюдаемых строго квантованы:
  \begin{align}
    \bigl(B_1(3),\ B_2(3)\bigr) &= \bigl(36 = K_2(\mathcal{C}_{13}),\ 72\bigr), \qquad
    \bigl(\mu_{\mathrm{vac}}(3),\ K_P(3)^2\bigr) \;=\; (1,\ 4\,096), \label{eq:mono-quant-obs-12} \\
    \hat{n}_E(\Phi) &= \lfloor|\Phi|/41\rfloor \in \{0,\dots,6\}, \qquad
    q_e \in \{-1, 0, +1\}, \label{eq:mono-quant-obs-34} \\
    (d_1 r)_f &= 512\,w_f, \qquad
    T_{\mathrm{spinor}} \;=\; 2\,\phaseRingCount(3) \;=\; 1\,024 \leftrightarrow 4\pi, \label{eq:mono-quant-obs-56} \\
    (q_u,\ q_d) &= \left(+\frac{2}{3},\ -\frac{1}{3}\right), \qquad
    D(\mathbf{x}) \in [8\,192,\ 833\times 8\,192], \qquad
    N_{\mathrm{pack}}(K) = K^3. \label{eq:mono-quant-obs-78}
  \end{align}
\end{enumerate}
\end{theorem}

\begin{proof}
В пункте~(1) при $Q = 1$ и полном периоде кольца $\Phi = 512$ средний фазовый шаг на $12$ соседях равен $\lfloor 512/12 \rfloor = 42 > 41 = \Delta\varphi_{\mathrm{disc}}(3)$, поэтому при $R = 0$ условие подпорогового удержания нарушается. На первой сфере $R_{\mathrm{dress}} = 1$ квадраты расстояний $|u - v_0|_2^2$ от $v_0 = (0,1,1)$ по $12$ вершинам $O_2(3)$ принимают значения $0$ ($1$ вершина $v_0$), $2$ ($4$ смежные вершины), $4$ ($4$ вершины), $6$ ($2$ вершины) и $8$ ($1$ антиподальная вершина $-v_0$), что вместе с центром $(0,0,0)$ даёт разбиение $1 + 1 + 4 + (4 + 2 + 1) = 1 + 1 + 4 + 7 = 13$.

В пункте~(2) деление с остатком на $K_P(3) = 64$ задаёт биекцию $\mathbb{Z}_{512} \leftrightarrow \mathbb{Z}_8 \times \mathbb{Z}_{64}$ ($512 = 8 \times 64$). Для спинорного периода~\eqref{eq:mono-quant-obs-56} сдвиг на полуоборот Паули $\Phi_{\mathrm{Pauli}}(3) = 256$ инвертирует знак дискретной экспоненты $\mathcal{E}_3(\varphi + 256) = -\mathcal{E}_3(\varphi)$, а повторный сдвиг $256 + 256 = 512 \equiv 0 \pmod{512}$ возвращает исходное состояние; тем самым полный двулистный спинорный цикл ($4\pi$) на двухслойной ячейке Верле составляет $2 \times 512 = 1\,024$ такта (в точном соответствии с периодом возврата параболы Верле $2\phaseRingCount(3) = 1\,024$). Для партонных зарядов~\eqref{eq:mono-quant-obs-78} на барионном треугольнике $K_3$ из $d = 3$ вершин с двумя спиновыми ориентациями Паули ($\uparrow, \downarrow$) распределение целочисленного числа обмотки $w \in \mathbb{Z}$ по $3$ вершинам в спиновой пропорции $2 : (-1)$ определяет партонные доли $q_u = +2/3$ и $q_d = -1/3$, суммирующиеся в строго целые заряды всех барионных замкнутых треугольников:
\begin{align*}
  Q(p = uud) &= \frac{2}{3} + \frac{2}{3} - \frac{1}{3} \;=\; +1, \qquad
  Q(n = udd) \;=\; \frac{2}{3} - \frac{1}{3} - \frac{1}{3} \;=\; 0, \\
  Q(\Delta^{++} = uuu) &= +2, \qquad
  Q(\Delta^- = ddd) \;=\; -1.
\end{align*}
\end{proof}

\begin{theorem}[Этаж~2. Полный вихревой атлас Стандартной модели на $\mathcal{C}_{13}$ в естественных безразмерных координатах ДКМ: три поколения лептонов, закон Койде, матрицы смешивания PMNS и CKM, шесть ароматов кварков, электрослабые и хиггсовский бозоны, магнитные моменты нуклонов, каналы слабых распадов и безразмерная лестница до атома водорода]
\label{thm:mono-complete-sm-particle-atlas}
В единой естественной шкале ячейки~\eqref{eq:mono-natural-scale-ladder} ($dh = 1,\ dt = 1,\ s_0 = 1/2,\ E_0 = s_0/dt = 1/2$) на $13$-узловом кликовом комплексе $\mathcal{C}_{13}$ с геометрическим отношением сфер кубооктаэдра $\kappa_{\mathrm{geom}}^2 = R_{\mathrm{in}}^2 / R_{\mathrm{out}}^2 = 1/2$, квантом упаковки одного ребра $\delta_{\mathrm{pack}} = \kappa_{\mathrm{geom}}^2 / |O_2(3)| = 1/24$, дискретным фазовым отношением порога дискретности к кольцу $r_{\mathcal{M}} = \Delta\varphi_{\mathrm{disc}}(3)/\phaseRingCount(3) = 41/512$ и безразмерной электрослабой полкой $v / E_0 = r_{\mathcal{M}}^{-1/2}\alpha_{\mathrm{fs}}^{|{\triangle}(3)|} = (512/41)^{1/2}\alpha_{\mathrm{fs}}^8$ (\texttt{elementary\_\allowbreak particles\_\allowbreak electron\_\allowbreak rotor\_\allowbreak clique\_\allowbreak complex\_\allowbreak and\_\allowbreak nucleon\_\allowbreak moments\_\allowbreak operator}, \texttt{complete\_\allowbreak particle\_\allowbreak spectrum\_\allowbreak quarks\_\allowbreak bosons\_\allowbreak hadrons\_\allowbreak and\_\allowbreak cosmology\_\allowbreak operator}):
\begin{enumerate}[label=\textbf{(\arabic*)},leftmargin=*]
\item \textbf{Три поколения заряженных лептонов $(e, \mu, \tau)$, точный инвариант Койде $Q_{\mathrm{Koide}} = 2/3$ и матрица смешивания нейтрино PMNS:}
  безразмерное отношение массы одетого электрона к бозону Хиггса равно $m_e / m_H = \alpha_{\mathrm{fs}}^2 / N_\varphi(3) = \alpha_{\mathrm{fs}}^2 / 13$. На $d = 3$ пространственных стеках поколений $g \in \{0, 1, 2\}$ с начальным спинорным сдвигом $\delta_0 = 2/\mathrm{mod\_bits}(3) = 2/9$ амплитуды лептонов $\sqrt{m_g} \propto 1 + \kappa_{\mathrm{geom}}^{-1}\cos(2\pi g/3 + 2/9)$ тождественно удовлетворяют закону Койде $Q_{\mathrm{Koide}}(3) = (d - 1)/d = 2/3$ ($m_\mu / m_e \approx 206{,}7703$, $m_\tau / m_e \approx 3477{,}47$), а доли слоёв разбиения ротора электрона $(1, 4, 7)$ в $N_\varphi(3) = 13$ и в третьей сфере $S_3(3) = 92$ определяют три угла смешивания нейтрино PMNS:
  \begin{align}
    \sin^2\theta_{12}^{\mathrm{PMNS}} &= \frac{d_{\mathrm{same}}(V)}{N_\varphi(3)} \;=\; \frac{4}{13}, \label{eq:mono-pmns-angles} \\
    \sin^2\theta_{23}^{\mathrm{PMNS}} &= \frac{d_{\mathrm{out}}(V)}{N_\varphi(3)} \;=\; \frac{7}{13}, \label{eq:mono-pmns-23} \\
    \sin^2\theta_{13}^{\mathrm{PMNS}} &= \frac{1_{\mathrm{core}} + d_{\mathrm{in}}(V)}{S_3(3)} \;=\; \frac{2}{92} \;=\; \frac{1}{46}. \label{eq:mono-pmns-13}
  \end{align}
\item \textbf{Три ортогональные плоскости цвета на $K_3$, шесть кварковых ступеней масс $(u, d, s, c, b, t)$ и матрица CKM из разложения $\widehat{M}(3) = 4^2 + 9^2 = 97$:}
  в каждом радиальном барионном треугольнике $K_3 = \{(0,0,0), (0,1,1), (1,0,1)\}$ три его ребра лежат в трёх взаимно перпендикулярных координатных плоскостях $(yz, xz, xy)$, образуя $N_c = d = 3$ цветовых канала. Безразмерные отношения масс кварков и параметры Вольфенштейна матрицы CKM равны:
  \begin{align}
    \frac{m_u}{m_e} &= d_{\mathrm{in}}(\square) = 4, \qquad
    \frac{m_d}{m_e} = \mathrm{mod\_bits}(3) = 9, \qquad
    \left(\frac{m_u}{m_e}\right)^2 + \left(\frac{m_d}{m_e}\right)^2 = 97 = \widehat{M}(3), \label{eq:mono-quarks-ud} \\
    \frac{m_p}{m_s} &= S_1(3) - 2 = 10, \qquad
    \frac{m_c}{m_p} = \frac{d_{\mathrm{in}}(\square)}{d} = \frac{4}{3}, \label{eq:mono-quarks-sc} \\
    \frac{m_b}{\alpha_{\mathrm{fs}}\,v} &= \frac{d_{\mathrm{out}}(V)}{d} = \frac{7}{3}, \qquad
    \left(\frac{m_t^{(\mathrm{bare})}}{v}\right)^2 = \frac{1}{2}, \label{eq:mono-quarks-scbt} \\
    \lambda_{\mathrm{CKM}} &= \frac{d_{\mathrm{out}}(E)\,d_{\mathrm{out}}(V)}{|O_2(3)|\,N_\varphi(3)} = \frac{35}{156}, \qquad
    A_{\mathrm{CKM}} = \frac{10}{12} = \frac{5}{6}, \qquad
    \bar{\rho} + i\bar{\eta} = \frac{4 + 9i}{26}. \label{eq:mono-ckm-params}
  \end{align}
\item \textbf{Магнитные моменты нуклонов, безразмерное отношение $(m_n - m_p)/m_e = 81/32$, электрослабый угол Вайнберга (древесный и на масштабе $M_Z$), каналы слабых распадов и безразмерная энергия связи атома водорода:}
  потеря одного кванта упаковки $\delta_{\mathrm{pack}} = 1/24$ на каждой из $d = 3$ вершин $K_3$ даёт изоскалярную сумму $1 - 3/24 = 7/8$, магнитные моменты протона и нейтрона, безразмерное отношение разности масс нейтрона и протона к массе электрона, аксиальную константу $g_A$, полные числа открытых каналов распада мюона и нейтрона (и средний канальный вес на один луч звезды $w_{\mathrm{ray}}(n)$), древесный и однопетлевой угол Вайнберга и безразмерное отношение энергии связи атома водорода к $E_0$ на $11$-й ступени иерархии $\alpha_{\mathrm{fs}}^{11}$:
  \begin{align}
    \frac{\mu_p}{\mu_N} &= 3 - 5\delta_{\mathrm{pack}} = \frac{67}{24}, \qquad
    \frac{\mu_n}{\mu_N} = -2 + 2\delta_{\mathrm{pack}} = -\frac{23}{12}, \qquad
    \frac{\mu_p + \mu_n}{\mu_N} = 1 - 3\delta_{\mathrm{pack}} = \frac{7}{8}, \label{eq:mono-nucleon-moments} \\
    \frac{m_n - m_p}{m_e} &= \frac{2\,|\mathcal{E}_{\mathrm{Heis}}(3)|}{K_P(3)} = \frac{81}{32} = 2{,}53125, \qquad
    g_A = \frac{d_{\mathrm{out}}(E)}{d_{\mathrm{in}}(\square)} + \frac{1}{d\,N_\varphi(3)} = \frac{5}{4} + \frac{1}{39} = \frac{199}{156}, \label{eq:mono-weak-decays} \\
    N_{\mathrm{open}}(\mu) &= |O_2(3)|\cdot\bigl(\phaseRingCount(3) - N_\varphi(3)\bigr) = 12 \times 499 = 5\,988, \label{eq:mono-muon-channels} \\
    N_{\mathrm{open}}^{(\mathrm{tot})}(n) &= N_{\mathrm{HB}}(3)\cdot\bigl(|O_2(3)| - 1\bigr) = 63 \times 11 = 693, \qquad
    w_{\mathrm{ray}}(n) = \frac{N_{\mathrm{open}}^{(\mathrm{tot})}(n)}{|O_2(3)|} = \frac{693}{12} = \frac{231}{4}, \label{eq:mono-neutron-channels} \\
    \sin^2\theta_W^{(0)} &= \frac{d}{N_\varphi(3)} = \frac{3}{13}, \qquad
    \sin^2\theta_W(M_Z) = \frac{d}{N_\varphi(3)} + \frac{\kappa_{\mathrm{geom}}^2\,\alpha_{\mathrm{fs}}}{|{\triangle}(3)|} = \frac{3}{13} + \frac{\alpha_{\mathrm{fs}}}{16}, \label{eq:mono-weinberg-dressed} \\
    \frac{m_H^{(\mathrm{bare})}}{v} &= \frac{1}{2}, \qquad
    \frac{\mathrm{BE}_{\mathrm{H}}}{E_0} = \frac{\alpha_{\mathrm{fs}}^2}{2}\,\frac{m_e}{E_0} = \frac{\alpha_{\mathrm{fs}}^{11}}{4\,N_\varphi(3)\,r_{\mathcal{M}}^{1/2}}. \label{eq:mono-bosons-hydrogen}
  \end{align}
\end{enumerate}
\end{theorem}

\begin{proof}
Проведём вывод каждого из пунктов~(1)--(3).
\begin{enumerate}[label=\textbf{\arabic*.},leftmargin=*]
\item Для трёх поколений лептонов $f(g) = 1 + \sqrt{2}\cos(2\pi g/3 + \delta_0)$ ($g \in \{0, 1, 2\}$) тригонометрические суммы по корням третьей степени из единицы $\sum_{g=0}^2 \cos(2\pi g/3 + \delta_0) = 0$ и $\sum_{g=0}^2 \cos^2(2\pi g/3 + \delta_0) = 3/2$ дают $\sum_{g=0}^2 f(g) = 3$ и $\sum_{g=0}^2 f(g)^2 = 3 + 2(3/2) = 6$, откуда отношение Койде равно $Q_{\mathrm{Koide}}(3) = 6 / 3^2 = 2/3$ тождественно при любом сдвиге фазы $\delta_0$. В разбиении $13$-узлового ротора электрона~\eqref{eq:mono-electron-13-partition} на $1_{\mathrm{core}} + 1 + 4 + 7 = 13$ доля $4$ смежных вершин квадратной грани даёт солнечный угол $\sin^2\theta_{12} = 4/13$, доля $7$ внешних несмежных вершин --- атмосферный угол $\sin^2\theta_{23} = 7/13$, а отношение $2$ осевых узлов $(0,0,0)$ и $v_0$ к мощности третьей координационной сферы $S_3(3) = 92$ --- реакторный угол $\sin^2\theta_{13} = 2/92 = 1/46$.
\item Для кварков первого поколения гауссов кулоновский модуль~\eqref{eq:mono-mhat-97} $\widehat{M}(3) = d_{\mathrm{in}}(\square)^2 + \mathrm{mod\_bits}(3)^2 = 4^2 + 9^2 = 97$ однозначно раскладывается в сумму двух квадратов с компонентами $m_u/m_e = d_{\mathrm{in}}(\square) = 4$ и $m_d/m_e = \mathrm{mod\_bits}(3) = 9$. На втором и третьем этажах отношение массы протона к странному кварку определяется главным коэффициентом Архимеда первой сферы $S_1(3) - (c_0/c)^2 = 12 - 2 = 10$, отношение очарованного кварка к протону --- отношением числа вершин квадратной грани к числу цветовых плоскостей $d_{\mathrm{in}}(\square)/d = 4/3$, отношение прелестного кварка к электрослабому масштабу $\alpha_{\mathrm{fs}}v$ --- отношением числа открытых октантов вершины к числу цветов $d_{\mathrm{out}}(V)/d = 7/3$, а голая константа Юкавы топ-кварка на геометрической проекции ребра равна $\kappa_{\mathrm{geom}}^2 = 1/2$. В матрице CKM угол Кабиббо равен отношению произведения внешних индексов ребра и вершины $d_{\mathrm{out}}(E)d_{\mathrm{out}}(V) = 5\times 7 = 35$ к полному числу направленных пар ядра $|O_2(3)|N_\varphi(3) = 12\times 13 = 156$, параметр $A_{\mathrm{CKM}} = (S_1(3)-2)/|O_2(3)| = 10/12 = 5/6$, а комплексная вершина апекса на удвоенном ядре $2 N_\varphi(3) = 26$ равна $(m_u + i m_d)/(2 N_\varphi m_e) = (4 + 9i)/26$.
\item На барионном треугольнике $K_3$ потеря одного кванта упаковки ребра $\delta_{\mathrm{pack}} = \kappa_{\mathrm{geom}}^2 / |O_2(3)| = 1/24$ на каждой из $d = 3$ вершин даёт изоскалярный спиновый момент $1 - 3\delta_{\mathrm{pack}} = 7/8$; вычитание $d_{\mathrm{out}}(E) = 5$ квантов упаковки из $d = 3$ магнетонов протона даёт $\mu_p/\mu_N = 3 - 5/24 = 67/24$, а добавление $(c_0/c)^2 = 2$ квантов упаковки к $-2$ магнетонам нейтрона даёт $\mu_n/\mu_N = -2 + 2/24 = -23/12$. Отношение удвоенного числа состояний беспотерьного коридора $2|\mathcal{E}_{\mathrm{Heis}}(3)| = 2\times 81 = 162$ к мантиссе $K_P(3) = 64$ даёт безразмерное отношение разности масс нуклонов к массе электрона $(m_n - m_p)/m_e = 162/64 = 81/32 = 2{,}53125$; отношение внешнего индекса ребра к квадратному входу $d_{\mathrm{out}}(E)/d_{\mathrm{in}}(\square) = 5/4$ плюс поправка одного цветового канала ядра $1/(d N_\varphi(3)) = 1/39$ даёт $g_A = 5/4 + 1/39 = 199/156$. В трёхчастичном распаде мюона $\mu^- \to e^- + \bar{\nu}_e + \nu_\mu$ каждый из $|O_2(3)| = 12$ выходных лучей кубооктаэдра располагает $\phaseRingCount(3) - N_\varphi(3) = 512 - 13 = 499$ свободными фазовыми состояниями кольца вне $13$-узлового ядра конечного электрона, что даёт $N_{\mathrm{open}}(\mu) = 12 \times 499 = 5\,988$ открытых каналов; в $\beta$-распаде нейтрона $N_{\mathrm{HB}}(3) = K_P(3) - 1 = 63$ ненулевых подуровня мантиссы по $|O_2(3)| - 1 = 11$ свободным лучам оболочки образуют $N_{\mathrm{open}}^{(\mathrm{tot})}(n) = 63 \times 11 = 693$ открытых мантиссово-лучевых канала со средним весом $w_{\mathrm{ray}}(n) = 693 / 12 = 231/4$ в расчёте на один из $|O_2(3)| = 12$ лучей кубооктаэдра. Наконец, для электрослабого угла Вайнберга~\eqref{eq:mono-weinberg-dressed} отношение числа пространственных осей $d = 3$ генераторов $\mathfrak{su}(2)_L$ к полному числу узлов причинного ядра $N_\varphi(3) = 13$ даёт древесное значение $\sin^2\theta_W^{(0)} = 3/13$, а однопетлевой электромагнитный сдвиг вдоль диагонального ребра ($\kappa_{\mathrm{geom}}^2 = 1/2$), распределённый по $|{\triangle}(3)| = 8$ октантным ступеням электрослабой иерархии ($\alpha_{\mathrm{fs}}^8$), прибавляет $\kappa_{\mathrm{geom}}^2\alpha_{\mathrm{fs}} / |{\triangle}(3)| = (1/2)\alpha_{\mathrm{fs}} / 8 = \alpha_{\mathrm{fs}} / 16$, приводя к $\sin^2\theta_W(M_Z) = 3/13 + \alpha_{\mathrm{fs}}/16 \approx 0{,}231225$. Подстановка $m_H = v/2$, $m_e = \alpha_{\mathrm{fs}}^2 m_H / N_\varphi(3)$ и ридберговского множителя $\alpha_{\mathrm{fs}}^2 / 2$ при $v / E_0 = r_{\mathcal{M}}^{-1/2}\alpha_{\mathrm{fs}}^8$ даёт безразмерное отношение $\mathrm{BE}_{\mathrm{H}} / E_0 = \alpha_{\mathrm{fs}}^{11} / (4 N_\varphi(3) r_{\mathcal{M}}^{1/2})$~\eqref{eq:mono-bosons-hydrogen}.
\end{enumerate}
\end{proof}

\paragraph{Единственный проектор перевода из естественной шкалы ДКМ в единицы СИ ($\pi_{\mathcal{M}\to\mathrm{SI}}$).}
Все теоремы монографии на носителе~$\mathcal{M}$ формулируются исключительно в естественных полярных координатах $(\mathbf{x}, t;\ \mu, \varphi_c)$ и безразмерных отношениях на единственной естественной механической шкале ячейки~\eqref{eq:mono-natural-scale-ladder} $(dh = 1,\ dt = 1,\ s_0 = 1/2,\ E_0 = s_0/dt)$. Перевод безразмерных инвариантов в конвенциональные единицы СИ (`м`, `с`, `эВ`, `нм`, `мкм`) осуществляется \textbf{в одном-единственном месте} --- при проекции $\pi_{\mathcal{M}\to\mathrm{SI}}$ через три фиксированных эталона СИ (постоянную Планка $\hbar = 2 s_0$, макроскопическую скорость света $c = \kappa_{\mathrm{geom}} dh / dt$ при $dh \leftrightarrow \ell_{\mathrm{P}}$ и $dt \leftrightarrow h_T = \kappa_{\mathrm{geom}} t_{\mathrm{P}}$, и энергетический масштаб $E_0 = s_0 / dt$, где дискретному фазовому фактору $r_{\mathcal{M}} = \Delta\varphi_{\mathrm{disc}}(3)/\phaseRingCount(3) = 41/512 \approx 0{,}080078$ в гладком сферическом пределе $\mathcal{T}$-слоя отвечает телесный фактор $1/(4\pi) \approx 0{,}079577$ с боровским радиусом $a_0 = N_{a0}\,dh = \hbar / (m_e c \alpha_{\mathrm{fs}}) \approx 0{,}529177\times 10^{-10}\text{ м}$ и ридберговской энергией $\mathrm{BE}_{\mathrm{H}} = \frac{1}{2}\alpha_{\mathrm{fs}}^2 m_e c^2 \approx 13{,}6057\text{ эВ}$). В частности, на $\mathcal{M}$-слое приведённый (радиальный) шаг фотона $\bar{\lambda}(E) := s_0 c / E$ в единицах боровского масштаба $\bar{\lambda}_0 := a_0\alpha_{\mathrm{fs}}^{-1} \approx 7{,}25163\text{ нм}$ является чисто алгебраическим и не содержит~$\pi$ ($\widehat{\Lambda}_{\mathrm{Ly}} := \bar{\lambda}_{\mathrm{Ly}} / (a_0\alpha_{\mathrm{fs}}^{-1}) = d - 1 = 2$), тогда как переход к циклической длине волны полного фазового оборота в лабораторном $\mathcal{T}$-спектрографе умножает $\bar{\lambda}$ на $2\pi$:
\begin{equation}
  \label{eq:mono-si-wavelength-bridge}
  \lambda_{\mathrm{SI}}(\widehat{\Lambda}) \;=\; 2\pi\,\bar{\lambda} \;=\; \widehat{\Lambda}\cdot\bigl(2\pi\,a_0\,\alpha_{\mathrm{fs}}^{-1}\bigr),
  \qquad
  2\pi\,a_0\,\alpha_{\mathrm{fs}}^{-1} \;=\; \frac{1}{2 R_\infty} \;\approx\; 45{,}56335\text{ нм}
\end{equation}
(так что лаймановскому порогу $\widehat{\Lambda}_{\mathrm{Ly}} = 2$ отвечает ридберговский предел $\lambda_\infty = R_\infty^{-1} = 4\pi\,a_0\,\alpha_{\mathrm{fs}}^{-1} \approx 91{,}1267\text{ нм}$, а трём мягким водородным окнам $\widehat{\Lambda}^{(+1,0,-1)} \in \{63,\ 126,\ 441\}$ --- инфракрасные длины волн $\lambda_{\mathrm{SI}}^{(+1,0,-1)} \approx \{2{,}870,\ 5{,}741,\ 20{,}09\}\text{ мкм}$).

\begin{table}[htbp]
\centering
\caption{Сводный атлас безразмерных теоретических значений ДКМ на $O_2(3)\times\mathbb{Z}_{512}$ в сопоставлении с экспериментальными данными CODATA-2022 и PDG (при нуле подгоночных параметров).}
\label{tab:mono-sm-codata-pdg-atlas}
\resizebox{\textwidth}{!}{%
\begin{tabular}{@{}llllr@{}}
\toprule
\textbf{Безразмерная наблюдаемая} & \textbf{Структура на $O_2(3)\times\mathbb{Z}_{512}$} & \textbf{Замкнутая формула ДКМ} & \textbf{Теория ДКМ} & \textbf{CODATA / PDG} \\
\midrule
Постоянная $\alpha_{\mathrm{fs}}^{-1}$ & $4$ симплекса $K_1\dots K_4$, $\widehat{M}=97$ & $(A_{\mathrm{den}} + B_{\mathrm{den}}\sqrt{2}) / (A_{\mathrm{num}} + B_{\mathrm{num}}\sqrt{2})$ & $137{,}035999187$ & $137{,}035999177$ \\
Угол Вайнберга $\sin^2\theta_W(M_Z)$ & $d=3$ оси в $N_\varphi(3)=13$, 1-петля & $3/13 + \alpha_{\mathrm{fs}}/16$ & $0{,}231225$ & $0{,}23122(4)$ \\
Константа $\alpha_s(M_Z)$ & $|{\triangle}(3)|=8$ граней $K_3$ & Бегунок от $\frac{3}{2}\cdot\frac{41}{512}$ & $0{,}11788$ & $0{,}1179(9)$ \\
Койде $Q_{\mathrm{Koide}}$ и $m_{\mu,\tau}/m_e$ & $3$ стека, $\delta_0=2/9$ & $Q_{\mathrm{Koide}} = 2/3$ & $206{,}77;\ 3477{,}5$ & $206{,}77;\ 3477{,}2$ \\
Углы PMNS $\sin^2\theta_{12,23,13}$ & Слои $(1,4,7)$ в $13$ и $92$ & $4/13,\quad 7/13,\quad 1/46$ & $0{,}308;\ 0{,}538$ & $0{,}307;\ 0{,}546$ \\
Параметры CKM $(\lambda,\ A)$ & $\widehat{M}=4^2+9^2=97$ на $156$ & $\lambda = 35/156,\quad A = 5/6$ & $0{,}224;\ 0{,}833$ & $0{,}225;\ 0{,}826$ \\
Апекс CKM $(\bar{\rho},\ \bar{\eta})$ & $(4 + 9i)$ на $2 N_\varphi = 26$ & $\bar{\rho} = 2/13,\quad \bar{\eta} = 9/26$ & $0{,}154;\ 0{,}346$ & $0{,}159;\ 0{,}348$ \\
Разность масс $(m_n - m_p)/m_e$ & Коридор $2|\mathcal{E}_{\mathrm{Heis}}|=162$ к $K_P=64$ & $162 / 64 = 81 / 32$ & $2{,}53125$ & $2{,}53099$ \\
Момент протона $\mu_p/\mu_N$ & Треугольник $K_3$, $\delta_{\mathrm{pack}}=1/24$ & $3 - 5/24 = 67/24$ & $+2{,}7917$ & $+2{,}7928$ \\
Момент нейтрона $\mu_n/\mu_N$ & Треугольник $K_3$ ($Q=0$) & $-2 + 2/24 = -23/12$ & $-1{,}9167$ & $-1{,}9130$ \\
Изоскаляр $(\mu_p+\mu_n)/\mu_N$ & Вычет $3\delta_{\mathrm{pack}} = 1/8$ & $1 - 1/8 = 7/8$ & $0{,}8750$ & $0{,}8798$ \\
Аксиальная константа $g_A$ & $d_{\mathrm{out}}(E)/d_{\mathrm{in}}(\square) + 1/39$ & $5/4 + 1/39 = 199/156$ & $1{,}27564$ & $1{,}2754(13)$ \\
Первичная доля гелия-4 $Y_p$ & $208 = 4\times 45 + 28$ на $t=2$ & $45 / (208 - 26) = 45/182$ & $0{,}24725$ & $0{,}245(3)$ \\
Гало к барионам $\Omega_{\mathrm{halo}}/\Omega_b$ & Куб грани $2^4$ к $d=3$ осям & $16 / 3$ ($\Omega_{\mathrm{tot}}/\Omega_b = 19/3$) & $5{,}333$ ($6{,}333$) & $5{,}36(5)$ ($6{,}36$) \\
\bottomrule
\end{tabular}%
}
\end{table}

\section{Гравитация как преломление на импедансе ворот без сингулярности \texorpdfstring{$1/0$}{1/0}, ряд сжатия Гаусса--Грина, кавитированные гиперъядра и фазы нейтронной пасты}

\begin{theorem}[Рефракционная гравитация без сингулярности $1/0$, ряд сжатия Гаусса--Грина, кавитированные гиперъядра, фазы нейтронной пасты и гало Вороного $16/3$ без тёмной материи]
\label{thm:mono-gravity-pasta-halo}
На решётке носителя при $d = 3$ локальный фазовый импеданс ячейки $D(\mathbf{x})$ удовлетворяет следующим законам (\texttt{gravity\_\allowbreak denominator\_\allowbreak strictly\_\allowbreak positive\_\allowbreak floor\_\allowbreak 8192}, \texttt{shell\_\allowbreak compression\_\allowbreak gauss\_\allowbreak flux\_\allowbreak limit\_\allowbreak and\_\allowbreak pauli\_\allowbreak rebound\_\allowbreak operator}, \texttt{gauss\_\allowbreak green\_\allowbreak compression\_\allowbreak series\_\allowbreak denominators\_\allowbreak and\_\allowbreak rational\_\allowbreak bounds}, \texttt{post\_\allowbreak island\_\allowbreak cavitated\_\allowbreak hypernuclei\_\allowbreak neutron\_\allowbreak pasta\_\allowbreak and\_\allowbreak nonsingular\_\allowbreak black\_\allowbreak hole\_\allowbreak core}, \texttt{no\_\allowbreak dark\_\allowbreak matter\_\allowbreak and\_\allowbreak vacuum\_\allowbreak holonomy\_\allowbreak halo}):
\begin{enumerate}[label=\textbf{(\arabic*)},leftmargin=*]
\item \textbf{Отсутствие сингулярности $1/0$ во всех конфигурациях поля и потолок противофазного стазиса $D_{\max} = 833\,\varepsilon(3)$:}
  для сырого знаменателя ворот $D_{\mathrm{raw}}(\mathbf{x})$ с вакуумным полом $\varepsilon(3) = K_P(3) Q_\square(3) = 64 \times 128 = 8\,192 = 2^{13}$
  \begin{equation}
    \label{eq:mono-gravity-denom-raw}
    D_{\mathrm{raw}}(\mathbf{x}) \;=\; \varepsilon(3) + Q_\square(3)\,\mu(\mathbf{x})^2 + \mu(\mathbf{x}) \sum_{\mathbf{e} \in O_2(3)} \mu(\mathbf{x}-\mathbf{e})\,\mathcal{C}_3\bigl(\Delta\varphi_{\mathbf{e}}(\mathbf{x})\bigr)
  \end{equation}
  защищённый оператор фазового импеданса ворот (\texttt{gateDenomOp 3})
  \begin{equation}
    \label{eq:mono-gravity-denom}
    D(\mathbf{x}) \;=\; \mathbf{1}_{\{D_{\mathrm{raw}}(\mathbf{x}) \ge 1\}}\,D_{\mathrm{raw}}(\mathbf{x}) + \mathbf{1}_{\{D_{\mathrm{raw}}(\mathbf{x}) \le 0\}}\,\varepsilon(3) \;\ge\; 1
  \end{equation}
  строго положителен при любых целочисленных значениях амплитуд и голономий ($D(\mathbf{x}) \ge 1$, причём на горизонте противофазного запирания Паули $D_{\mathrm{raw}}(\mathbf{x}) \le 0$ принимает точное значение вакуумного пола $D(\mathbf{x}) = \varepsilon(3) = 8\,192 > 0$). В однородном кластере максимальной плотности $\mu(\mathbf{x}) = \mu(\mathbf{x}-\mathbf{e}) = K_P(3) = 64$ при синфазном насыщении $\mathcal{C}_3(0) = +128$ по всем $12$ лучам импеданс достигает конечного целочисленного потолка противофазного стазиса:
  \begin{equation}
    \label{eq:mono-stasis-833}
    D_{\max}(3) \;=\; \bigl(1 + N_\varphi(3)\cdot K_P(3)\bigr)\,\varepsilon(3) \;=\; (1 + 13\times 64)\times 8\,192 \;=\; 833 \times 8\,192 \;=\; 6\,823\,936,
  \end{equation}
  где $D_{\max}/\varepsilon = 833 = 7^2 \times 17 = d_{\mathrm{out}}(V)^2\,(N_\varphi + d_{\mathrm{in}}(\square))$, радиальный ток на противофазном горизонте обращается в нуль ($\mathcal{S}_3(256) = 0$), а побитовая обратимость шага Верле сохраняет всю информацию без потерь (устраняя гравитационную сингулярность $1/0$ и информационный парадокс Хокинга).
\item \textbf{Ряд сжатия Гаусса--Грина, кавитированные гиперъядра и три топологические фазы нейтронной пасты:}
  число межсферных рёбер Гаусса между сферами $k-1$ и $k$ равно $E(k-1\leftrightarrow k) = 12(k^3 - (k-1)^3)$ ($\sum_{k=1}^K E(k-1\leftrightarrow k) = 12 K^3$), порождая семь канонических знаменателей ряда сжатия Гаусса--Грина $m^3 - (m-1)^3 \in \{1, 7, 19, 37, 61, 91, 127 = Q_\square - 1\}$ при $m = 1,\dots,7$. Выше порога кавитации центрального узла $Z_{\mathrm{Dirac}}^{(k=0)} = 3\times 41 = 123$ ($u_0 \equiv 0$) полая первая сфера $S_1$ устойчива до $Z_{\mathrm{hollow},S_1} = 10\times 3^3 = 270 = \Phi_{\mathrm{Pauli}} + Z(\mathrm{Si})$, а под внешним гравитационным давлением кора нейтронной звезды проходит три топологические фазы «ядерной пасты» ($\chi_{\mathrm{2D}} = 2$ «ньокки» $\to \chi_{\mathrm{2D}} = 0$ «спагетти» $\to \chi_{\mathrm{2D}} = 41$ $6$-слойная «лазанья» с $V = 42, E = 147$).
\item \textbf{Вакуумное гало Вороного $\Omega_{\mathrm{halo}}/\Omega_b = 16/3$ без частиц тёмной материи ($D_{\mathrm{dark}} \equiv 0$):}
  отношение полного $16$-состоянийного булева куба грани $2^{d_{\mathrm{in}}(\square)} = 16$ к $d = 3$ пространственным проекциям даёт коэффициент вакуумного гало Вороного и полную долю материи:
  \begin{align}
    \frac{\Omega_{\mathrm{halo}}}{\Omega_b} &= \frac{2^{d_{\mathrm{in}}(\square)}}{d} \;=\; \frac{16}{3} \;\approx\; 5{,}333, \label{eq:mono-voronoi-halo-16-3} \\
    \frac{\Omega_{\mathrm{halo}} + \Omega_b}{\Omega_b} &= \frac{16 + 3}{3} \;=\; \frac{19}{3} \;=\; \frac{N_\varphi(3) + |\square(3)|}{d} \;\approx\; 6{,}333. \label{eq:mono-voronoi-halo-19-3}
  \end{align}
\end{enumerate}
\end{theorem}

\begin{proof}
В пункте~(1) при произвольных амплитудах соседей $\mu(\mathbf{x}-\mathbf{e}) \in [1, 64]$ и противофазных разностях $\Delta\varphi_{\mathbf{e}}$ сумма $\mathcal{Z}_{\mathrm{real}}(\mathbf{x})$ может быть отрицательной и превосходить по модулю $128\mu(\mathbf{x})^2 + 8\,192$ (как показано в доказательстве Теоремы~\ref{thm:mono-first-vortex-birth-iv263}, в узле $(-2,-1,-1)$ при $t=1$ сырой знаменатель равен $D_{\mathrm{raw}} = -113\,826 < 0$). В определении оператора $\mathrm{clampDenomOp}_3$~\eqref{eq:mono-gravity-denom} при $D_{\mathrm{raw}} \le 0$ целочисленное усечение в $\mathbb{N}$ даёт $\texttt{denomRaw.toNat} = 0$, индикатор $\texttt{isZeroNat}(0) = 1$ выбирает первую ветвь линейной склейки $\varepsilon(3) = 8\,192 > 0$, а при $D_{\mathrm{raw}} = k + 1 \ge 1$ индикатор $\texttt{isZeroNat}(k+1) = 0$ сохраняет $D_{\mathrm{raw}} \ge 1$; тем самым $D(\mathbf{x}) > 0$ для всех $(\mu, \mathcal{Z}_{\mathrm{real}}) \in \mathbb{Z}^2$ (\texttt{gravity\_denominator\_strictly\_positive\_floor\_8192}). При синфазном же насыщении однородного кластера $\mu(\mathbf{x}) = \mu(\mathbf{x}-\mathbf{e}) = K_P(3) = 64$ и $\mathcal{C}_3(0) = +128$ все $13$ узлов причинного ядра $B_1(3) = \{\mathbf{0}\}\cup O_2(3)$ дают положительный вклад $128 \times 64^2 = 64 \times 8\,192 = K_P(3)\,\varepsilon(3)$, откуда
\[
  D_{\max}(3) \;=\; \varepsilon(3) + 64\,\varepsilon(3) + 12 \times 64\,\varepsilon(3) \;=\; (1 + 13 \times 64)\,\varepsilon(3) \;=\; 833\,\varepsilon(3) \;=\; 6\,823\,936.
\]
В пункте~(2) телескопическое суммирование $12\sum_{k=1}^K(k^3 - (k-1)^3) = 12 K^3$ и подстановка $m = 1,\dots,7$ в $m^3 - (m-1)^3$ дают $\{1, 7, 19, 37, 61, 91, 127\}$, а на $6$-слойной колонке нейтронной лазаньи $E - V - K_P = 147 - 42 - 64 = 41$. В пункте~(3) $2^4 / 3 = 16/3$ и $(16 + 3)/3 = 19/3 = (13 + 6)/3$.
\end{proof}

